\section{Conclusion}

In this paper, we introduced the Regime-Enhanced Intelligent Granger Network (REIGN), a principled and scalable framework for causal discovery in nonstationary multivariate time series. REIGN directly addresses the critical gap in existing literature: classical and state-of-the-art time-series causal discovery algorithms uniformly assume stationarity, rendering them structurally unreliable in the presence of regime shifts, distributional breaks, and time-varying causal mechanisms that characterize real-world financial, physiological, and climate systems.

REIGN resolves this challenge through a unified four-stage pipeline. First, MICE-based preprocessing ensures data quality without compromising inter-variable covariance structure. Second, the PELT changepoint detection algorithm partitions the time series into locally stationary regimes with $\mathcal{O}(T)$ complexity, enabling each downstream model to specialize on a homogeneous distribution. Third, a zero-shot Large Language Model querying mechanism synthesizes domain-specific causal priors that constrain the vast DAG search space without requiring labeled data or human expert annotation. Fourth, a per-regime Message-Passing GNN optimizes a continuous adjacency matrix through an Augmented Lagrangian objective, where the LLM prior regularizes the loss landscape and acyclicity is enforced differentiably via the NOTEARS-style matrix exponential constraint.

Our empirical evaluation on synthetic nonstationary VAR data demonstrates that REIGN achieves definitive state-of-the-art performance, outperforming PCMCI+ by over 11 percentage points in Optimal F1 score and dominating across AUROC and AUPR metrics. The ablation study reveals that both the regime-aware ensemble mechanism and the dense GNN initialization are individually critical architectural contributors, with the ensemble providing the largest single marginal improvement.

\subsection{Limitations}

The primary limitation of REIGN in its current form is computational: the $\mathcal{O}(N^3)$ cost of the matrix exponential DAG constraint is prohibitive for very large graphs ($N > 100$). Furthermore, the PELT segmentation operates independently of the causal model, meaning the detected regimes may not be optimally aligned with the true causal structural breaks. Future work will explore joint optimization of the changepoint locations and the causal model parameters.

\subsection{Future Directions}

Several promising directions emerge from this work. First, replacing the heuristic PELT-based segmentation with a differentiable Hidden Markov Model (HMM) would enable end-to-end gradient-based optimization of both the regime boundaries and the causal models simultaneously. Second, incorporating Chebyshev polynomial approximations of the DAG constraint would reduce the per-step complexity from $\mathcal{O}(N^3)$ to $\mathcal{O}(N^2)$, making REIGN tractable for datasets with hundreds of variables. Third, extending the LLM prior mechanism to incorporate retrieval-augmented generation (RAG) over domain-specific literature would significantly enrich the quality of the zero-shot topological scaffolds provided to the GNN. Finally, validating REIGN on real-world financial datasets — such as multi-asset returns across quantitative regime shifts identified by central bank policy announcements — would establish its industrial applicability and open new directions for interpretable causal econometrics.
